\subsection{Le condizioni minime di esistenza}\label{sec:c-e}

Uno stream è istanziato a partire da tre
parametri indispensabili — \texttt{stream\_id}, \texttt{onset}, \texttt{sample} — che
ne costituiscono le condizioni minime di esistenza, come mostrato nel
frammento di \texttt{identity.yml} (Listato~\ref{lst:c-e}).
Il dizionario dello stream ha quindi una notazione differenziale: ciò che non
è scritto è intervento nullo, e ogni scarto udibile
dalla sorgente è imputabile a un parametro esplicito.

\begin{figure}[H]
    \begin{framedfloat}
    \lstinputlisting[
      linerange={3-7},
      language=yaml,
      basicstyle=\ttfamily\fontsize{7.5}{9}\selectfont,
      label=lst:c-e,
      caption={Le tre chiavi obbligatorie
      (\texttt{identity.yml}): in assenza di una \texttt{duration} dichiarata,
      lo stream dura quanto il file audio che legge.
      Infatti lo stato corrente deve coincidere con la trasformazione identica della
      sorgente (con finestra di Hann e fattore di overlap 2: configurazione che garantisce una risintesi
      a guadagno unitario e inalterata\protect\footnote{Finestrare un grano ne modula l'ampiezza,
aggiungendo allo spettro bande laterali~\cite{Roads2001}. Nella configurazione minima questi contributi si elidono nella
sovrapposizione dei grani: la differenza RMS gain-matched fra sorgente e
resa non supera i $-74$~dB, sotto la soglia udibile.}).
      Il file audio \texttt{voices.wav} verrà utilizzato
      in tutti gli esempi successivi.}
    ]{examples/identity/identity.yml}
    \end{framedfloat}
\end{figure}

